\documentclass[5p,times,twocolumn,authoryear]{elsarticle}

% Paquetes esenciales para idioma y matemáticas
\usepackage[spanish]{babel}
\usepackage[utf8]{inputenc}
\usepackage{amsmath}
\usepackage{amssymb}
\usepackage{booktabs}
\usepackage{hyperref}
\usepackage{graphicx}
\graphicspath{{../Panel/}{../Panel/figures/}}
\usepackage{longtable}
\usepackage{array}

\begin{document}

\begin{frontmatter}

\title{Fragilidad Bancaria y Riesgos del Crecimiento a la Baja: Cómo la Interacción entre JLoss y GaR Determina el Spread Soberano en Economías en Desarrollo}

%% Información del autor
\author[1]{Mauricio Valenzuela Corvalán}
\affiliation[1]{organization={Magíster en Economía Aplicada, Universidad de Chile},
            city={Santiago},
            country={Chile}}

\begin{abstract}
Este bosquejo de artículo evalúa empíricamente cómo la interacción no lineal entre la fragilidad sistémica del sector bancario y los riesgos de la cola izquierda del crecimiento del producto determina el riesgo de crédito soberano en economías en desarrollo. La fragilidad bancaria se mide mediante la métrica \textit{Joint Loss} (JLoss) (Chari et al., 2024), un indicador semiparamétrico a nivel país de la pérdida conjunta esperada del sistema bancario condicional a un evento sistémico, construido a partir de probabilidades de incumplimiento de tipo Merton (1974) agregadas vía aproximación de punto de silla (Martin et al., 2001). El riesgo a la baja del crecimiento se captura mediante \textit{Growth-at-Risk} (GaR), estimado por regresión cuantílica de panel con efectos fijos (Koenker, 2004) sobre un índice de condiciones financieras, siguiendo el marco de \textit{vulnerable growth} de Adrian, Boyarchenko y Giannone (2019) y la implementación de CEMLA (Ossandon Busch et al., 2022). Sobre el núcleo de cinco economías latinoamericanas (Brasil, Chile, Colombia, México y Perú), el coeficiente de interacción $JLoss\times GaR$ es negativo y significativo tanto en la ventana original (2010Q1--2022Q2, $N=248$; $\hat\theta=-0{,}352$ a $-0{,}363$, $t\approx-3{,}4$ a $-4{,}0$) como en una ventana ampliada hacia atrás mediante un GaR reentrenado con mayor historia (2007Q4--2022Q2, $N=293$; $\hat\theta=-0{,}305$ a $-0{,}391$, $p<0{,}03$ en todas las especificaciones), y es robusto a la incorporación —como índices agregados— de los controles domésticos y globales del paper de referencia (deuda/PIB, margen de utilidad bancaria, tasa del Tesoro a 10 años y \textit{spread} on/off-the-run). Una extensión exploratoria a 11 economías emergentes con EMBI real (BCRP y tablas del \textit{Global Financial Stability Report} del FMI) no reproduce el mismo patrón de forma limpia fuera del núcleo latinoamericano: en la ventana común 2010Q1--2014Q4, el subpanel de solo-LatAm no rechaza $\theta=0$ ($\hat\theta=-0{,}074$, $p=0{,}56$), mientras que el panel de 11 países en la misma ventana entrega un coeficiente positivo y marginalmente significativo ($\hat\theta=+0{,}626$, $p=0{,}05$), un hallazgo que se reporta con honestidad como evidencia de heterogeneidad no modelada —no como refutación— y que delimita el alcance actual de la generalización del mecanismo. Un modelo de umbral de panel (Hansen, 2000) sobre el núcleo LatAm corrobora la no linealidad de la interacción. La evidencia respalda la hipótesis de complementariedad —fragilidad bancaria y riesgo de cola del crecimiento se amplifican mutuamente en su transmisión al spread soberano— de manera robusta dentro de América Latina, y motiva la ampliación de la cobertura de datos como condición para extender la conclusión al resto de la muestra de economías emergentes.
\end{abstract}

\begin{keyword}
estabilidad financiera \sep fragilidad bancaria \sep growth-at-risk \sep spread soberano \sep doom loop \sep economías emergentes
\JEL G01 \sep G21 \sep F34 \sep E44 \sep C21
\end{keyword}

\end{frontmatter}

\section{Introducción}

\begin{itemize}
    \item \textbf{Motivación y hechos estilizados.} Durante las crisis financieras recientes, los spreads soberanos de las economías emergentes se han ampliado de manera más pronunciada que los de las economías avanzadas. Se argumenta que parte de esta ampliación es atribuible a la fragilidad del sector bancario doméstico, que eleva la deuda pública contingente y los riesgos de balance, y a su interacción con un debilitamiento de las perspectivas de crecimiento.
    \item \textbf{Brecha en la literatura.} La literatura ha tratado por separado (i) el nexo banca–soberano (fragilidad bancaria $\rightarrow$ riesgo soberano) y (ii) el nexo finanzas–crecimiento (condiciones financieras $\rightarrow$ cola izquierda del PIB). No existe evidencia sistemática que articule ambos canales como determinantes conjuntos del spread soberano, ni que caracterice la naturaleza no lineal de su interacción.
    \item \textbf{Pregunta de investigación.} ¿Amplifica la coincidencia de alta fragilidad bancaria ($JLoss$) y elevado riesgo a la baja del crecimiento ($GaR$) el riesgo soberano por encima de la suma de sus efectos individuales?
    \item \textbf{Aporte (triple contribución, en la estructura de Chari et al., 2024):}
    \begin{enumerate}
        \item Integra dos métricas estado-del-arte —una de fragilidad bancaria sistémica (JLoss) y una de vulnerabilidad macroeconómica (GaR)— en un único marco de determinación del spread soberano para economías emergentes.
        \item Caracteriza la \textbf{no linealidad} del nexo banca–soberano–crecimiento mediante un término de interacción y un modelo de umbral, contrastando si la amplificación es dependiente de régimen.
        \item Provee evidencia con un \textit{pipeline} reproducible de extracción de balances bancarios y de estimación de GaR pseudo–tiempo real, triangulando el resultado central en tres bases de datos —núcleo LatAm en dos ventanas temporales y una extensión exploratoria a 11 economías emergentes— para delimitar con precisión dónde el mecanismo está establecido y dónde constituye todavía una hipótesis abierta.
    \end{enumerate}
    \item \textbf{Hoja de ruta.} La Sección 2 desarrolla el marco teórico. La Sección 3 sintetiza el estado del arte. La Sección 4 detalla la metodología de construcción de JLoss y GaR y la estrategia econométrica. La Sección 5 describe los datos y las tres bases utilizadas. La Sección 6 reporta los resultados, incluida la triangulación entre bases. La Sección 7 discute las implicancias y concluye.
\end{itemize}

\section{Marco Teórico}

El argumento central de este trabajo se sostiene sobre tres pilares teóricos que la literatura ha desarrollado de forma mayoritariamente independiente. El primero es el bucle de retroalimentación soberano-bancario (\textit{doom loop}), que articula la transmisión bidireccional de fragilidad entre los balances del sistema bancario y los del soberano. El segundo es la interpretación de la fragilidad bancaria sistémica como un pasivo contingente del soberano, que provee el fundamento microeconómico de la métrica $JLoss$. El tercero es la determinación del riesgo a la baja del crecimiento por las condiciones financieras —el programa de investigación de \textit{vulnerable growth}—, que sustenta la métrica $GaR$.

La contribución conceptual del paper no reside en ninguno de estos pilares por separado, sino en la tesis de que operan de manera \textbf{complementaria} y no aditiva: la fragilidad bancaria y la vulnerabilidad del crecimiento se amplifican mutuamente en su transmisión al riesgo soberano. Esta sección desarrolla cada pilar y formaliza la predicción de amplificación supra-aditiva que orienta la estrategia empírica de las Secciones 4 a 6.

\subsection{El nexo soberano-bancario: del bucle de retroalimentación al abrazo mortal}

La piedra angular teórica es el mecanismo de retroalimentación entre la solvencia bancaria y la solvencia soberana formalizado por Farhi y Tirole (2018). En su modelo de tres fechas, el soberano emite deuda cuyo precio refleja la capacidad fiscal esperada, mientras los bancos domésticos gestionan su liquidez manteniendo bonos soberanos —domésticos y extranjeros— como reserva de valor.

Ante un choque adverso en la fecha intermedia, la solvencia del soberano se deteriora por dos vías: un efecto \textbf{directo}, en tanto el choque contrae la actividad económica y, con ella, los ingresos fiscales; y un efecto \textbf{indirecto}, que opera a través de los balances bancarios. La caída inicial del precio de la deuda pública erosiona el patrimonio neto de los bancos y deprime su inversión; si el soberano valora suficientemente esa inversión, rescata al sistema bancario. El punto decisivo es que el rescate se financia emitiendo deuda adicional, lo que reduce todavía más su precio: por cada dólar de rescate, el acervo de deuda pública debe aumentar en más de un dólar.

Formalmente, este mecanismo admite una representación de punto fijo para el precio de la deuda soberana doméstica en la fecha intermedia, $p_1(s) = 1 - F\!\left(B_1(s)\mid s\right)$, donde $B_1(s)$ es el acervo de deuda estado-contingente y $F(\cdot)$ la distribución de la capacidad fiscal. La sensibilidad de $p_1(s)$ al estado $s$ adopta la forma de un \textbf{multiplicador de amplificación}: una caída del precio incrementa el rescate requerido y el volumen de bonos a colocar, lo que vuelve a deprimir el precio, \textit{ad infinitum}. La magnitud de este multiplicador es creciente en el grado de \textit{home bias} —la proporción de deuda soberana en manos de bancos domésticos—.

La implicación que este trabajo explota es que la amplificación es \textbf{intrínsecamente no lineal}: el deterioro soberano no es la suma de un componente bancario y uno fiscal independientes, sino su producto a través de un multiplicador que se intensifica con la fragilidad del balance bancario y con la severidad del estado adverso.

\subsection{La fragilidad bancaria como pasivo contingente del soberano: el canal $JLoss$}

La métrica de fragilidad bancaria empleada en este trabajo, $JLoss$, se interpreta como el costo directo esperado de rescatar al conjunto del sistema bancario ante un evento sistémico, y por ello entra naturalmente en la valoración del bono soberano (Chari et al., 2024). Esta interpretación operacionaliza el ``\textit{bailout put}'' de Farhi y Tirole (2018) y la dilución de Acharya et al. (2014): si el mercado anticipa que un soberano deberá asumir el costo de recapitalizar su banca en un escenario de tensión, ese pasivo contingente debe estar incorporado \textit{ex ante} en la prima soberana.

\subsection{El riesgo a la baja del crecimiento: el canal $GaR$}

El tercer pilar proviene del programa de \textit{vulnerable growth} de Adrian, Boyarchenko y Giannone (2019), cuya tesis es que las condiciones financieras predicen de manera \textbf{asimétrica} la distribución condicional del crecimiento futuro del producto: la dependencia respecto de las condiciones financieras es marcadamente más fuerte en los cuantiles inferiores —la cola izquierda— que en los superiores, que permanecen comparativamente estables en el tiempo. Sus hallazgos asociados —la correlación negativa entre la media y la volatilidad condicionales del crecimiento, y la ``paradoja de la volatilidad'' según la cual períodos de baja volatilidad anteceden a resultados de crecimiento negativos (Brunnermeier y Sannikov, 2014; Adrian y Boyarchenko, 2012)— implican que el objeto relevante para la evaluación del riesgo no es el crecimiento esperado, sino la severidad de su escenario adverso plausible.

La traducción de este canal al riesgo soberano opera a través de la dinámica de sostenibilidad de la deuda. La ley de movimiento del ratio de deuda sobre producto, $\dot{d} = (r - g)\,d - pb$ —donde $d$ es la deuda/PIB, $r$ la tasa de interés real, $g$ la tasa de crecimiento y $pb$ el balance primario—, muestra que un sesgo a la baja en la distribución de crecimiento deteriora la trayectoria esperada de sostenibilidad y, sobre todo, la capacidad fiscal precisamente en los estados de cola donde el incumplimiento se torna probable. Dado que el costo del default posee un perfil de pago de tipo opción —discreto y concentrado en el peor escenario—, lo determinante para la prima soberana no es el primer momento de la distribución de crecimiento, sino su cola izquierda. Es esta no linealidad la que conecta el canal de crecimiento con el canal soberano y la que hace insuficiente controlar por el crecimiento medio: se requiere una medida del riesgo de cola, esto es, $GaR$.

Operativamente, $GaR$ se define como el cuantil de cola izquierda de la distribución condicional del crecimiento (con la convención de que valores más negativos indican mayor riesgo a la baja), estimado mediante regresión cuantílica (Koenker y Bassett, 1978; Koenker, 2004) sobre un índice de condiciones financieras. El concepto, nacido como extensión macroeconómica del \textit{value-at-risk} (Prasad et al., 2019), se implementa en este trabajo siguiendo la plataforma de código abierto de CEMLA (Ossandon Busch et al., 2022), que organiza los factores de riesgo del crecimiento en tres bloques —condiciones macrofinancieras locales, un índice de condiciones financieras y condiciones internacionales ortogonalizadas— y caracteriza la densidad condicional completa del crecimiento en lugar de únicamente su media. El puente conceptual entre condiciones financieras laxas, acumulación de vulnerabilidad y materialización del riesgo de cola lo ofrece Stein (2012), al concebir la política monetaria como instrumento de regulación de la estabilidad financiera.

\subsection{Síntesis: complementariedad y amplificación supra-aditiva}

Los tres pilares convergen en una predicción común que constituye la hipótesis teórica central del trabajo: \textbf{el efecto de la fragilidad bancaria sobre el riesgo soberano no es independiente del estado del crecimiento, y viceversa}. La coincidencia de alta fragilidad bancaria y elevado riesgo a la baja del crecimiento amplifica el spread soberano por encima de la suma de los efectos individuales de cada dimensión.

La justificación de esta complementariedad se desprende de la articulación de los tres pilares:

\begin{itemize}
    \item \textbf{Del mecanismo de Farhi y Tirole (2018):} el multiplicador de amplificación es creciente tanto en la fragilidad del balance bancario como en la severidad del estado adverso $s$. Identificando $GaR$ con la cola izquierda del crecimiento —es decir, con los estados adversos $s$—, el efecto cruzado del multiplicador implica que la pérdida de precio del bono soberano ante un incremento de la fragilidad es mayor en estados de bajo crecimiento. La amplificación es, por construcción, un fenómeno de interacción.
    \item \textbf{Del canal de dilución de Acharya et al. (2014):} la sensibilidad del riesgo soberano a las condiciones bancarias escala con el grado de deterioro pre-existente, y el factor macroeconómico común que vincula banca y soberano es el crecimiento. La interacción fragilidad $\times$ estado macro es, por tanto, una predicción directa del modelo.
    \item \textbf{De la complementariedad de Gennaioli et al. (2014):} el costo del incumplimiento soberano crece con la exposición bancaria; un escenario de bajo crecimiento que eleva la probabilidad de default, combinado con un sistema bancario frágil, magnifica el costo total de la cesación de pagos y, con él, la prima soberana.
\end{itemize}

Estos argumentos motivan el objeto empírico central del trabajo: el término de interacción $JLoss_{i,t}\times GaR_{i,t}$ y la predicción de \textbf{complementariedad (super-aditividad)}. Bajo la convención adoptada —$GaR$ medido en niveles como el cuantil de cola izquierda del crecimiento, de modo que valores \textbf{más negativos} señalan mayor riesgo a la baja—, el efecto marginal de la fragilidad bancaria sobre el spread debe intensificarse a medida que el riesgo de cola se agrava. Formalmente, la complementariedad se expresa como

\[ \frac{\partial^{2}\, Spread_{i,t}}{\partial\, JLoss_{i,t}\;\partial\, GaR_{i,t}} \;<\; 0, \]

esto es, el efecto marginal de la fragilidad bancaria sobre el spread es \textbf{decreciente en $GaR$}: a menor $GaR$ (mayor severidad de la cola), mayor es la sensibilidad del spread a $JLoss$. En la especificación lineal con interacción, $\partial\, Spread/\partial\, JLoss = \beta_{1} + \theta\,GaR$, por lo que la hipótesis de amplificación se traduce en un coeficiente de interacción \textbf{negativo} ($\theta < 0$): cuando $GaR$ toma valores marcadamente negativos, el término $\theta\,GaR$ es positivo y eleva el efecto de la fragilidad. Una representación de régimen mediante un modelo de umbral (Hansen, 2000) formaliza la versión no lineal de esta hipótesis: existe un umbral de vulnerabilidad $\lambda$ por debajo del cual (en el régimen de $GaR$ suficientemente adverso) la sensibilidad del spread a la fragilidad bancaria es estructuralmente mayor.

Esta formulación delimita con precisión la contribución del trabajo frente a la literatura previa. Chari et al. (2024) modelan la interacción de $JLoss$ con factores financieros \textbf{globales} —el VIX, la tasa del Tesoro estadounidense, los spreads de alto rendimiento—, mostrando que los países con banca más frágil son más sensibles a los choques exógenos del ciclo financiero global. El presente trabajo desplaza el foco hacia la interacción de $JLoss$ con una métrica de vulnerabilidad \textbf{doméstica y endógena del crecimiento} ($GaR$), cerrando así el triángulo banca-crecimiento-soberano que la literatura ha analizado únicamente por pares. Las hipótesis contrastables que se derivan del marco son:

\begin{itemize}
    \item \textbf{H1 (canal bancario).} La fragilidad bancaria sistémica se asocia positivamente con el spread soberano: $\beta_{1} > 0$.
    \item \textbf{H2 (canal de crecimiento).} Un mayor riesgo a la baja del crecimiento —un $GaR$ más negativo— se asocia con un mayor spread soberano; en niveles, el signo esperado es $\beta_{2} < 0$.
    \item \textbf{H3 (complementariedad — hipótesis central).} La interacción entre fragilidad bancaria y riesgo a la baja del crecimiento amplifica el spread soberano de manera supra-aditiva, con coeficiente de interacción $\theta < 0$; de forma equivalente, existe un régimen de alta vulnerabilidad ($GaR$ por debajo de un umbral $\lambda$) en el que la sensibilidad del spread a sus determinantes es estructuralmente superior.
\end{itemize}

% ---

\section{Estado del Arte}

Esta sección sitúa el trabajo en la intersección de cuatro literaturas que han evolucionado de manera predominantemente paralela: la de los determinantes del spread soberano en economías emergentes, la de las métricas de riesgo sistémico y fragilidad bancaria, la del nexo banca-soberano, y la de la frontera entre condiciones financieras y crecimiento del producto (\textit{growth-at-risk}). La revisión persigue un objetivo argumentativo preciso: mostrar que, pese a la madurez de cada literatura por separado, ninguna ha articulado de forma conjunta y no lineal la fragilidad bancaria sistémica y el riesgo a la baja del crecimiento como codeterminantes del riesgo soberano —que es, justamente, la brecha que este trabajo cubre.

\subsection{Determinantes del spread soberano en economías emergentes}

La literatura clásica sobre el riesgo soberano de las economías emergentes ha oscilado entre dos polos explicativos. En un extremo, los \textbf{fundamentos macroeconómicos} domésticos —razón deuda/PIB, crecimiento, términos de intercambio, reservas, calidad institucional— como determinantes del diferencial soberano (Edwards, 1984; Hilscher y Nosbusch, 2010). En el otro, los \textbf{factores globales} o de oferta de capital —el apetito por riesgo, la tasa libre de riesgo estadounidense, la liquidez internacional—, sintetizados en la distinción entre factores \textit{push} y \textit{pull} de la literatura de flujos de capital (Calvo, Leiderman y Reinhart, 1996). La evidencia ha tendido a favorecer la preeminencia de los factores globales: González-Rozada y Levy-Yeyati (2008) muestran que una fracción sustancial de la variación de los spreads de los mercados emergentes se explica por el apetito de riesgo global y los spreads de alto rendimiento de las economías avanzadas; Longstaff, Pan, Pedersen y Singleton (2011) documentan que el riesgo de crédito soberano es mayoritariamente sistemático y se halla más vinculado a los mercados financieros estadounidenses que a los fundamentos locales; y Remolona, Scatigna y Wu (2008) descomponen el spread en una pérdida esperada y una prima de riesgo, atribuyendo la dinámica de esta última a la aversión global al riesgo. La conexión entre el ciclo económico doméstico y el spread fue formalizada por Uribe y Yue (2006), mientras que la consolidación de un único factor financiero global que gobierna la covariación de los precios de los activos riesgosos a escala mundial —el \textit{Global Financial Cycle}— se debe a Miranda-Agrippino y Rey (2022). Cavallo y Valenzuela (2010) y Csonto e Ivaschenko (2013), por su parte, descomponen el papel relativo de fundamentos y factores globales en la determinación de los diferenciales. Este trabajo se inserta en esa tradición pero introduce un determinante doméstico hasta ahora ausente del núcleo de la literatura —el riesgo a la baja del crecimiento medido como cuantil de cola— e indaga su interacción con la fragilidad bancaria, en lugar de tratar a ambos como regresores aditivos independientes.

\subsection{Métricas de riesgo sistémico y fragilidad bancaria}

La crisis financiera global de 2008-09 catalizó una vasta literatura orientada a medir el riesgo sistémico del sector financiero. Pueden distinguirse tres familias metodológicas. La primera, basada en datos de mercado y en la teoría de carteras, incluye el \textit{Marginal Expected Shortfall} y el \textit{SRISK} —la insuficiencia de capital condicional del sistema ante un evento sistémico— de Acharya, Pedersen, Philippon y Richardson (2017) y Brownlees y Engle (2016), así como el \textit{CoVaR} y su variación $\Delta$CoVaR de Adrian y Brunnermeier (2016), que cuantifican el riesgo de cola del sistema condicional a la dificultad de una institución individual. La segunda familia, de naturaleza estructural y semiparamétrica, deriva de los modelos de Merton (1974) y agrega las pérdidas potenciales del sistema mediante técnicas de distribución de pérdidas de cartera: el \textit{Distress Insurance Premium} de Huang, Zhou y Zhu (2009) —la prima hipotética de un seguro contra pérdidas sistémicas— y las \textit{Banking Stability Measures} de Segoviano y Goodhart (2009) son sus exponentes más cercanos a la métrica empleada en este trabajo. La tercera familia, de corte econométrico, mide la interconexión y la propagación, como en los índices de conectividad de Billio, Getmansky, Lo y Pelizzon (2012); una panorámica de este instrumental se encuentra en Bisias, Flood, Lo y Valavanis (2012).

La métrica $JLoss$ pertenece a la segunda familia y, en rigor, es pariente directa del \textit{Distress Insurance Premium} de Huang et al. (2009): ambas construyen la distribución agregada de pérdidas del sistema a partir de probabilidades de incumplimiento de tipo Merton y de la dependencia respecto de un factor sistémico. Su especificidad —y su ventaja para el problema soberano— reside en tres rasgos: la agregación a \textbf{nivel de país} (la unidad pertinente para el riesgo soberano y no, como en buena parte de la literatura anglosajona, a nivel de institución o de sistema estadounidense), la \textbf{interpretación como costo directo de rescate} y, decisivamente, la \textbf{disponibilidad para una muestra amplia de economías emergentes}, donde la cobertura de SRISK, CoVaR y medidas afines es escasa o inexistente. La capacidad predictiva comparada de estas métricas sobre la actividad real fue evaluada por Giglio, Kelly y Pruitt (2016), cuyo conjunto de medidas sirve además como batería de controles de robustez en la literatura aplicada y en el presente trabajo.

\subsection{El nexo banca-soberano}

La interdependencia entre la solvencia bancaria y la soberana constituye el corazón teórico del paper y ha generado una literatura propia, revitalizada por la crisis de deuda europea. En el plano histórico-empírico, Reinhart y Rogoff (2009, 2011) documentan la recurrencia con que las crisis bancarias anteceden y desencadenan crisis soberanas, y Laeven y Valencia (2013, 2020) cuantifican los cuantiosos costos fiscales de las crisis bancarias sistémicas en su base de datos de referencia. En el plano teórico, las contribuciones seminales ya discutidas en el marco —Farhi y Tirole (2018), Acharya, Drechsler y Schnabl (2014), Gennaioli, Martin y Rossi (2014)— se complementan con los modelos de riesgo soberano e intermediación de Bocola (2016), el análisis de fragilidad bancaria en mercados financieramente integrados de Bolton y Jeanne (2011), el rol de las garantías gubernamentales en la retroalimentación de Leonello (2018) y la propuesta de activos seguros europeos (ESBies) de Brunnermeier et al. (2016) como dispositivo para romper el bucle.

En el plano empírico aplicado, Mody y Sandri (2012) muestran que bancos y soberanos quedaron "unidos por la cadera" durante la crisis europea; Dieckmann y Plank (2012) documentan la transferencia de riesgo del sector privado al público y su modulación por la importancia pre-crisis del sistema financiero doméstico; Kallestrup, Lando y Murgoci (2016) evidencian que los vínculos del sector financiero contienen información sobre el CDS soberano más allá de los fundamentos; y Fratzscher y Rieth (2019) analizan la interacción entre política monetaria, rescates bancarios y el nexo soberano-bancario. La dimensión de la ``\textit{carry trade}'' —bancos que cargan deuda soberana riesgosa de su propio país— fue documentada por Acharya y Steffen (2015). Esta literatura, sin embargo, es abrumadoramente europea y de frecuencia de crisis; su extensión sistemática a un panel amplio de economías emergentes, con una métrica de fragilidad construida homogéneamente para todos los países, es precisamente la línea que abren Chari et al. (2024) y que este trabajo profundiza al incorporar el canal de crecimiento.

\subsection{\textit{Growth-at-Risk} y la frontera finanzas-crecimiento}

La cuarta literatura, de aparición más reciente, vincula el estado de las condiciones financieras con la distribución condicional —y no sólo con la media— del crecimiento futuro del producto. Su contribución fundacional es Adrian, Boyarchenko y Giannone (2019), que establece la asimetría esencial del fenómeno: las condiciones financieras predicen con fuerza la cola izquierda del crecimiento, mientras la cola derecha permanece estable. El marco fue adoptado por los organismos multilaterales como instrumento de vigilancia macroprudencial —el \textit{Global Financial Stability Report} del FMI de 2017 y la aplicación de Prasad et al. (2019)— y trasladado a las economías latinoamericanas mediante la plataforma de código abierto de CEMLA (Ossandon Busch et al., 2022), base metodológica directa de la estimación de $GaR$ de este trabajo. Antecedentes conceptuales del nexo finanzas-cola del crecimiento se hallan en Allen, Bali y Tang (2012), que indagan si una medida agregada de riesgo de cola del sistema financiero (CATFIN) predice contracciones macroeconómicas, y en Giglio, Kelly y Pruitt (2016). El sustento teórico del canal proviene de la literatura de fijación de precios de activos por intermediarios financieros —Brunnermeier y Sannikov (2014), He y Krishnamurthy (2013)—, que racionaliza por qué las fricciones financieras generan no linealidades y riesgo de cola en la actividad real.

Desde el punto de vista econométrico, el método descansa en la regresión cuantílica de Koenker y Bassett (1978), sistematizada en Koenker (2005). Su extensión a datos de panel con efectos fijos —el estimador empleado en este trabajo— plantea desafíos específicos abordados por Koenker (2004), con penalización de los efectos individuales, y por las alternativas de Canay (2011) y Machado y Santos Silva (2019); la monotonicidad de la función cuantílica estimada se garantiza mediante el reordenamiento de Chernozhukov, Fernández-Val y Galichon (2010). La construcción del índice de condiciones financieras que alimenta la estimación se apoya en la tradición de Hatzius et al. (2010) y en indicadores compuestos de estrés sistémico como el CISS de Holló, Kremer y Lo Duca (2012). Es de notar que esta literatura ha mantenido su variable de interés en el crecimiento del producto; su conexión con el \textit{pricing} del riesgo soberano permanece, hasta donde alcanza esta revisión, esencialmente inexplorada.

\subsection{Brecha en la literatura y posicionamiento del trabajo}

La revisión precedente revela una compartimentación persistente. La literatura del nexo banca-soberano (3.3) ha establecido que la fragilidad bancaria eleva el riesgo soberano, pero ha tratado el crecimiento como un fundamento de fondo o como un factor de confusión a controlar —no como un canal de cola con dinámica propia—. La literatura de \textit{growth-at-risk} (3.4) ha caracterizado con precisión cómo las condiciones financieras deterioran la cola izquierda del crecimiento, pero ha detenido su análisis en el producto, sin proyectarlo sobre la prima soberana. Y la literatura de determinantes del spread (3.1) ha privilegiado los factores globales y los fundamentos en niveles, omitiendo las métricas de riesgo de cola del crecimiento. El trabajo más próximo, Chari et al. (2024), introduce la fragilidad bancaria sistémica ($JLoss$) en la determinación del spread soberano de economías emergentes y explora su interacción con factores financieros \textbf{globales}, pero no con una medida doméstica del riesgo a la baja del crecimiento.

La contribución de este artículo consiste en cerrar el triángulo banca-crecimiento-soberano. Frente a esa literatura fragmentada, el trabajo (i) emplea conjuntamente una métrica estructural de fragilidad bancaria sistémica ($JLoss$) y una métrica de riesgo de cola del crecimiento ($GaR$) como determinantes del spread soberano; (ii) contrasta la naturaleza \textbf{no lineal} de su interacción —la hipótesis de complementariedad supra-aditiva— mediante un término de interacción y un modelo de umbral (Hansen, 2000); y (iii) triangula el resultado en tres bases de datos que varían en cobertura temporal y transversal, delimitando explícitamente el alcance actual de la evidencia. En ese sentido, el aporte no es la mera adición de un regresor, sino la articulación teórica y empírica de dos canales que la literatura ha mantenido separados, reportada con el nivel de honestidad estadística que exige extender la conclusión más allá de la muestra donde fue establecida.

% ---

\section{Metodología}

La estrategia metodológica se articula en tres bloques. Los dos primeros describen la construcción de las variables explicativas de interés —la fragilidad bancaria sistémica $JLoss_{i,t}$ (Sección 4.1) y el riesgo a la baja del crecimiento $GaR_{i,t}$ (Sección 4.2)—, ambas estimadas a partir de modelos estructurales y semiparamétricos sobre microdatos bancarios y macrofinancieros. El tercer bloque (Sección 4.3) detalla la estrategia de identificación econométrica con la que se contrastan las hipótesis H1–H3. Las dos métricas se construyen de manera independiente y se consolidan en un panel a frecuencia trimestral mediante combinación por país y trimestre.

\subsection{Medición de la fragilidad bancaria: la métrica $JLoss$}

$JLoss_{i,t}$ es una métrica semiparamétrica, a nivel de país, de la pérdida conjunta esperada del sistema bancario condicional a un evento sistémico, expresada como porcentaje de la exposición total. Su construcción procede en cuatro etapas: (i) la estimación de las probabilidades individuales de incumplimiento de cada banco mediante un modelo estructural de tipo Merton; (ii) la condicionalización de esas probabilidades respecto de un factor sistémico país; (iii) la agregación de la distribución de pérdidas del sistema por el método de punto de silla; y (iv) el cálculo de las contribuciones marginales y la normalización que define la métrica. Se sigue la metodología de Chari et al. (2024).

\subsubsection{Probabilidades individuales de incumplimiento}

Para cada banco se estima la probabilidad de incumplimiento siguiendo la modificación de Kealhofer (2000) del modelo estructural de Merton (1974), en su variante KMV. A partir de los pasivos de corto y largo plazo ($L_{ST}$, $L_{LT}$), se define el \textbf{punto de incumplimiento}

\[ D^{*} = L_{ST} + \tfrac{1}{2}\,L_{LT}, \]

donde el factor $\tfrac{1}{2}$ sobre la deuda de largo plazo recoge la convención empírica de KMV. Dados el valor de mercado del patrimonio $E$ (capitalización bursátil), su volatilidad $\sigma_E$, la tasa libre de riesgo $r$ y el horizonte $T$, se resuelve numéricamente el sistema de dos ecuaciones no lineales que proyecta el valor de los activos $\hat V_A$ y su volatilidad implícita $\hat\sigma_A$:

\[ E = \hat V_A\,\Phi(d_1) - D^{*}\,e^{-rT}\,\Phi(d_2), \qquad \sigma_E\,E = \Phi(d_1)\,\hat\sigma_A\,\hat V_A, \]

con $d_1 = \dfrac{\ln(\hat V_A/D^{*}) + \left(r + \tfrac{1}{2}\hat\sigma_A^{2}\right)T}{\hat\sigma_A\sqrt{T}}$ y $d_2 = d_1 - \hat\sigma_A\sqrt{T}$, donde $\Phi(\cdot)$ es la función de distribución normal estándar. La volatilidad del patrimonio se estima por \textbf{varianza realizada} (Barndorff-Nielsen et al., 2002): para cada trimestre y banco, $\sigma_E = \sqrt{\sum_{t=1}^{Q} r_t^{2}}$, donde $r_t$ son los retornos logarítmicos diarios del patrimonio y $Q$ el número de días de transacción del trimestre. Con las soluciones $\hat V_A$ y $\hat\sigma_A$ se calcula la \textbf{distancia al incumplimiento} y la \textbf{frecuencia esperada de incumplimiento}:

\[ DD = \frac{\hat V_A - D^{*}}{\hat V_A\,\hat\sigma_A}, \qquad EDF = \Phi(-DD). \]

El valor $EDF$ de cada banco constituye su probabilidad incondicional de incumplimiento $p_i \equiv pdef_i$, insumo del método de agregación.

\subsubsection{Factor sistémico y probabilidades condicionales}

La hipótesis identificadora de la agregación es la \textbf{independencia condicional}: los riesgos bancarios son independientes una vez condicionados a un único factor sistémico país, que se toma como el retorno del índice bursátil de cada economía. Sea $\rho_i$ la correlación entre los retornos del patrimonio del banco $i$ y los del factor sistémico $V$. Siguiendo la estructura unifactorial de Vasicek (1997), la probabilidad de incumplimiento \textbf{condicional} al estado $V$ del factor es

\[ p_i(V) = \Phi\!\left(\frac{\Phi^{-1}(p_i) - \rho_i\,V}{\sqrt{1-\rho_i^{2}}}\right). \]

Esta expresión —implementada literalmente en la rutina \texttt{cond\_pd}— traslada la dependencia de cola al canal sistémico: cuando $V$ adopta valores adversos, las probabilidades condicionales de todos los bancos aumentan simultáneamente, generando la correlación de pérdidas que un agregado de riesgos idiosincráticos ignoraría.

\subsubsection{Agregación de pérdidas por el método de punto de silla}

La pérdida agregada del sistema se define como $P = \sum_{i=1}^{N} e_i\,\mathbf{1}_{D_i}$, donde $e_i$ es la exposición al incumplimiento del banco $i$ (proxy: pasivos totales) y $\mathbf{1}_{D_i}$ la función indicadora de incumplimiento. El método de punto de silla (Martin, Thompson y Browne, 2001) traslada el problema al espacio de la función generadora de momentos, donde admite tratamiento analítico. La función generadora de cumulantes (CGF) de la pérdida, condicional al factor $V$, es

\[ K(s\mid V) = \sum_{i=1}^{N} \ln\!\left(1 - p_i(V) + p_i(V)\,e^{s\,a_i}\right), \]

con $a_i$ la exposición normalizada del banco $i$. Sus derivadas primera y segunda —$K'(s)$ y $K''(s)$, implementadas como \texttt{get\_K1st} y \texttt{get\_K2nd}— proporcionan la media y la varianza de la distribución de pérdidas. Para un nivel de pérdida $x$, el \textbf{punto de silla} $\hat t$ resuelve $K'(\hat t) = x$. La probabilidad de cola se aproxima entonces, siguiendo Martin et al. (2001), por

\[ P(L > x \mid V) \approx 1 - \exp\!\left(K(\hat t\mid V) - \hat t\,x + \tfrac{1}{2}\hat t^{2} K''(\hat t)\right)\Phi\!\left(-|\hat t|\sqrt{K''(\hat t)}\right), \]

para $x$ por encima de la media (con la rama simétrica correspondiente en la cola inferior). La integración sobre la distribución del factor sistémico —supuesta $N(0,\sigma_V^{2})$— se resuelve por \textbf{cuadratura de Gauss–Hermite} con siete nodos:

\[ P(L > x) \approx \sum_{k=1}^{7} w_k\,P(L > x \mid V_k), \]

donde $\{V_k, w_k\}$ son los nodos y pesos de la cuadratura. La curva de probabilidad acumulada se evalúa sobre una grilla de 500 pasos entre las cotas de pérdida calibradas (1\% y 4,8\% de la exposición). El valor en riesgo al 99\% ($VaR_{99}$) se obtiene por interpolación lineal sobre esa curva, sin extrapolación (réplica de \texttt{LinealXY2}).

\subsubsection{Contribuciones marginales y definición de $JLoss$}

Siguiendo a Martin (2001), la contribución marginal del banco $n$ al riesgo del sistema, evaluada en el percentil 99, es

\[ C_n = e_n \cdot \frac{\sum_{k} \nu_k\,\tilde p_n(\hat t_k)}{\sum_{k}\nu_k}, \qquad \nu_k = w_k\,\frac{\exp\!\left(K(\hat t_k) - \hat t_k\,K'(\hat t_k)\right)}{\sqrt{K''(\hat t_k)}}, \qquad \tilde p_n(\hat t) = \frac{p_n\,e^{\hat t\,a_n}}{1 - p_n + p_n\,e^{\hat t\,a_n}}, \]

donde $\tilde p_n(\hat t)$ es la probabilidad exponencialmente inclinada (\textit{saddle-tilted}) del banco $n$. La suma de estas contribuciones en el percentil 99 constituye la \textbf{pérdida inesperada} ($UL$), mientras que la \textbf{pérdida esperada} es $EL = \sum_i e_i\,p_i$. La métrica de fragilidad se define finalmente como las pérdidas totales normalizadas por la exposición agregada:

\[ JLoss_{i,t} = \frac{EL_{i,t} + UL_{i,t}}{\sum_{j} EAD_{j,t}}\times 100, \]

donde la exposición al incumplimiento incorpora la pérdida dada el incumplimiento, $e_j = EAD_j\cdot LGD$, con $LGD = 45\%$ (BIS, 2006).

\subsubsection{Parametrización e implementación computacional}

El Cuadro 1 resume la parametrización, consistente con Chari et al. (2024): horizonte de un trimestre; $LGD = 45\%$; siete nodos de cuadratura; un único factor sistémico; cota inferior de pérdidas de 1,0\% y superior de 4,8\%; 500 pasos de discretización; percentil de 99\%; y factor sistémico $N(0,\,0{,}16\%)$. La totalidad del cómputo se reimplementó en Python validándolo contra el código original en MATLAB (\texttt{countrypd.m}, \texttt{loss\_distrib.m}, \texttt{loss\_contrib.m}, \texttt{get\_prob.m}, \texttt{LinealXY2.m}, \texttt{find\_saddlev1.m} y los módulos \texttt{get\_K*.m}), reproduciendo sus resultados. La búsqueda del punto de silla emplea el algoritmo de Brent (\texttt{brentq}) en lugar del Newton–Raphson original, por su mayor robustez ante curvaturas extremas de la CGF, e incorpora salvaguardas numéricas frente a desbordamiento en la evaluación de las exponenciales.

\vspace{1em}
\noindent \textbf{Cuadro 1 — Parámetros del método de punto de silla.} Horizonte: 1 trimestre $\cdot$ LGD: 45\% $\cdot$ Nodos de cuadratura: 7 $\cdot$ Factores sistémicos: 1 $\cdot$ Cota inferior de pérdidas: 0,010 $\cdot$ Cota superior: 0,048 $\cdot$ Pasos: 500 $\cdot$ Percentil: 0,99 $\cdot$ Factor sistémico: $N(0,\,0{,}16\%)$.
\vspace{1em}

\subsection{Medición del riesgo a la baja del crecimiento: la métrica $GaR$}

$GaR_{i,t}$ es el cuantil de cola izquierda (5\%) de la distribución condicional del crecimiento del producto un trimestre adelante, estimado a partir de un índice de condiciones financieras mediante regresión cuantílica de panel. Su construcción adapta el marco de \textit{vulnerable growth} de Adrian, Boyarchenko y Giannone (2019) y la plataforma de CEMLA (Ossandon Busch et al., 2022), con tres modificaciones metodológicas: estimación por programación lineal global, lectura directa del cuantil sin densidad kernel, y estimación pseudo–tiempo real.

\subsubsection{Índice de condiciones financieras (FCI)}

El FCI sintetiza por país un conjunto de indicadores de tensión financiera —de los mercados monetario, cambiario, accionario, de deuda soberana y bancario— junto con un bloque de tipo de cambio real efectivo. El procedimiento replica la cadena de rutinas del marco CEMLA y consta de tres pasos. Primero, cada indicador se \textbf{estandariza} mediante una transformación expansiva (\texttt{method\_b\_max}), consistente con la naturaleza pseudo–tiempo real de la estimación posterior. Segundo, se computan \textbf{correlaciones temporales por pares} entre los indicadores estandarizados mediante un suavizado exponencial (EWMA) con factor de decaimiento $\lambda = 0{,}85$. Tercero, el índice se agrega como una \textbf{forma cuadrática} sobre los indicadores estandarizados y la matriz de correlaciones EWMA $C_t$, siguiendo la lógica del indicador compuesto de estrés sistémico de Holló, Kremer y Lo Duca (2012):

\[ FCI_t = \sqrt{(w \circ s_t)'\,C_t\,(w \circ s_t)}, \]

donde $s_t$ es el vector de subíndices estandarizados y $w$ el vector de ponderaciones. La implementación se validó contra la salida de referencia de CEMLA (\texttt{FCI\_MEXICO.csv}), alcanzando una correlación de 0,9995 y un RMSE de 0,016 sobre 212 meses.

\subsubsection{Regresión cuantílica de panel con efectos fijos}

Para el cuantil $\tau$ y el horizonte $h$, la distribución condicional del crecimiento se caracteriza mediante la regresión cuantílica de panel

\[ Q\!\left(\Delta GDP_{i,t+h}\mid \cdot\,;\tau\right) = \alpha_h(\tau) + \beta_{1,h}(\tau)\,\Delta GDP_{i,t} + \beta_{2,h}(\tau)\,FCI_{i,t} + \beta_{3,h}(\tau)\,VIX_t + \mu_{i,h}(\tau), \]

con $h = 1$ trimestre y efectos fijos de país $\mu_i$ (Koenker, 2004). El VIX se ortogonaliza respecto de las condiciones locales. Los parámetros se obtienen minimizando la suma de pérdidas de comprobación,

\[ \min_{\beta,\,\mu}\;\sum_{i}\sum_{t}\rho_\tau\!\left(\Delta GDP_{i,t+h} - x_{i,t}'\beta - \mu_i\right), \qquad \rho_\tau(u) = u\,(\tau - \mathbf{1}\{u<0\}), \]

problema que se formula como un \textbf{único programa lineal} y se resuelve con el algoritmo HiGHS (vía \texttt{scipy.optimize.linprog}), superando la tolerancia del solver iterativo \texttt{rqpd} de R.

\subsubsection{De la función cuantílica a $GaR$}

La regresión se estima sobre una grilla de $n_\tau = 39$ cuantiles, diseñada de modo que $\tau = 0{,}05$ caiga exactamente sobre un nodo. Para garantizar la monotonicidad se aplica el \textbf{reordenamiento} de Chernozhukov, Fernández-Val y Galichon (2010):

\[ GaR_{i,t} \equiv Q\!\left(\Delta GDP_{i,t+1}\mid \cdot\,;\,\tau = 0{,}05\right). \]

\subsubsection{Estimación pseudo-tiempo real}

Para evitar el sesgo de anticipación (\textit{look-ahead bias}), la serie $\{GaR_{i,t}\}$ se construye con una \textbf{ventana expansiva}: en cada fecha $t$, los parámetros cuantílicos se reestiman utilizando exclusivamente la información disponible hasta $t$. La serie fue reentrenada durante esta entrega sobre un panel de 15 países con mayor historia disponible, lo que —para el núcleo LatAm— extiende la ventana estimable desde 2010Q1 hasta 2007Q4 (Sección 5.6) y habilita la extensión exploratoria a 11 economías emergentes con EMBI real (Sección 5.7).

\subsubsection{Robustez: corroboración con densidad skew-$t$}

Como control de robustez se ajusta la distribución \textit{skew-t} de Adrian, Boyarchenko y Giannone (2019) a los cuantiles condicionales estimados. La correlación entre esta serie y la obtenida por lectura directa del cuantil reordenado es de 0,9996.

\subsection{Estrategia de identificación econométrica}

\subsubsection{Especificación de referencia}

Siguiendo la lógica de identificación de Chari et al. (2024), la especificación de referencia es un modelo de panel con efectos fijos bidireccionales:

\[ EMBI_{i,t} = \alpha_i + \gamma_t + \beta_1\,JLoss_{i,t} + \beta_2\,GaR_{i,t} + \omega'X_{i,t} + \epsilon_{i,t}, \]

donde $\alpha_i$ son efectos fijos de país que absorben toda heterogeneidad invariante en el tiempo; $\gamma_t$ son efectos fijos de tiempo que absorben los choques globales comunes —de modo que, bajo $\gamma_t$, los factores financieros globales quedan absorbidos y no se incluyen explícitamente en las especificaciones bidireccionales—; y $X_{i,t}$ es el vector de controles.

\subsubsection{Interacción y complementariedad}

La hipótesis central (H3) se contrasta incorporando el término de interacción:

\[ EMBI_{i,t} = \alpha_i + \gamma_t + \beta_1\,JLoss_{i,t} + \beta_2\,GaR_{i,t} + \theta\,\bigl(JLoss_{i,t}\times GaR_{i,t}\bigr) + \omega'X_{i,t} + \epsilon_{i,t}. \]

Bajo la convención de signo adoptada, la complementariedad supra-aditiva predice $\theta < 0$. El efecto marginal de la fragilidad sobre el diferencial,

\[ \frac{\partial\,EMBI_{i,t}}{\partial\,JLoss_{i,t}} = \beta_1 + \theta\,GaR_{i,t}, \]

se reporta evaluado a distintos niveles de $GaR$ en un gráfico de efectos parciales análogo al de Chari et al. (2024).

\subsubsection{Modelo de umbral de panel}

La versión no lineal de la hipótesis de amplificación se contrasta mediante un modelo de umbral (Hansen, 1999, 2000):

\[ EMBI_{i,t} = \alpha_i + \gamma_t + \boldsymbol{\beta}^{(1)'} Z_{i,t}\,\mathbf{1}\{q_{i,t}\le\lambda\} + \boldsymbol{\beta}^{(2)'} Z_{i,t}\,\mathbf{1}\{q_{i,t}>\lambda\} + \omega'X_{i,t} + \epsilon_{i,t}, \]

donde $q_{i,t}=GaR_{i,t}$. El umbral $\lambda$ se estima por mínimos cuadrados concentrados tras la transformación \textit{within}; su significancia se evalúa mediante un contraste de razón de verosimilitud cuya distribución se aproxima por \textit{bootstrap} (Hansen, 2000).

\subsubsection{Construcción del panel, inferencia y controles agregados}

Las series de $JLoss$, $GaR$, EMBI y controles se consolidan por país y trimestre. La inferencia se basa en errores estándar de Driscoll–Kraay ($vcovSCC$), robustos simultáneamente a heterocedasticidad, autocorrelación serial y dependencia transversal —justificados por los diagnósticos de la Sección 6.8—, contrastados frente a \textit{clustering} por país, por tiempo y al estimador HAC de Arellano. Los seis controles adicionales replicados del paper de referencia (deuda/PIB, margen de utilidad bancaria, volatilidad cambiaria, tasa del Tesoro a 10 años, \textit{spread} de alto rendimiento y \textit{spread} on/off-the-run) se incorporan, cuando aplica, no como seis regresores sueltos sino como \textbf{dos índices agregados} —uno doméstico y uno global, Sección 6.11— dado que mezclar variables de variación puramente transversal-temporal con variables de variación estrictamente temporal bajo efectos fijos bidireccionales induce a que el componente global sea absorbido por los efectos de tiempo, atenuando artificialmente su poder identificador.

% ---

\section{Datos}

\subsection{Tres bases, un mismo objetivo}

La estrategia de construcción de datos contempla un universo objetivo de $\sim$19 economías emergentes a frecuencia trimestral. Dado que la disponibilidad de balances bancarios (insumo de $JLoss$), EMBI real y GaR difiere sustancialmente entre países, esta entrega reporta y triangula la evidencia sobre \textbf{tres bases} que representan distintos compromisos entre profundidad temporal, rigor de los controles y amplitud transversal:

\begin{itemize}
    \item \textbf{LatAm corta} — Brasil, Chile, Colombia, México y Perú; 2010Q1–2022Q2; $N=248$; GaR estimado sobre el panel original de 5 países. Es la base sobre la que las tres métricas están más plenamente validadas y sobre la que descansa el resultado central del trabajo.
    \item \textbf{LatAm larga} — mismos cinco países y mismo EMBI real (BCRP), pero con GaR reentrenado sobre un panel de 15 países con mayor historia de entrenamiento, lo que extiende la ventana estimable hasta 2007Q4; $N=293$. Permite auditar si "hacer más robusto" el GaR —incorporando más historia en su estimación pseudo-tiempo real— altera la conclusión central.
    \item \textbf{Extendida (11 países)} — agrega a los cinco países LatAm seis economías adicionales con EMBI \textbf{real}, no proxy (Bulgaria, China, Indonesia, Polonia, Sudáfrica y Turquía, recuperado de tablas del \textit{Global Financial Stability Report} del FMI, 2011–2015); $N=395$. Sin controles macroeconómicos domésticos para los seis países nuevos y con VIX global (CBOE) en vez del VIX/crecimiento de la plataforma CEMLA. Constituye la primera prueba exploratoria de generalización del mecanismo fuera de América Latina.
\end{itemize}

\subsection{Variable dependiente}
Spread soberano EMBI, en puntos básicos (especificación de nivel) y en logaritmo natural (especificación M4, elasticidad).

\subsection{Insumos de $JLoss$}
Hojas de balance bancarias estandarizadas, tasa de interés promedio, capitalización bursátil y volatilidad realizada de retornos accionarios, extraídos por país desde fuentes regulatorias nacionales (CMF, BCB IFData, SBS, CNBV, SBP, SARB, entre otras).

\subsection{Insumos de $GaR$}
Componentes del FCI por país (IPC, rendimiento soberano a 10 años, índice accionario, REER) y VIX como proxy de condiciones financieras globales.

\subsection{Variables de control}
Deuda/PIB, balance fiscal/PIB, reservas/PIB, cuenta corriente/PIB, inflación interanual y tipo de cambio real efectivo (controles domésticos, disponibles para el núcleo LatAm); VIX. Adicionalmente, replicando los controles del paper de referencia (Chari et al., 2024): margen de utilidad bancaria (\textit{profit margin}, reconstruido desde balances históricos 1999–2011), tasa del Tesoro de EE.UU. a 10 años, \textit{spread} de alto rendimiento (proxy ICE BofA US HY OAS) y \textit{spread} on/off-the-run construido con la curva cero-cupón de Gürkaynak-Sack-Wright de la Reserva Federal. La volatilidad cambiaria del paper de referencia queda pendiente de construcción.

\subsection{Estadística descriptiva: núcleo LatAm, corta vs. larga}

El Cuadro 2 compara los estadísticos descriptivos de las variables centrales entre la base corta y la larga. La ampliación de la ventana hacia 2007Q4 —que incorpora la crisis financiera global— eleva levemente la media y la dispersión del EMBI (de 213 a 221 pb en promedio) sin alterar cualitativamente ninguna otra variable, un primer indicio de que la ventana adicional no está distorsionando la muestra.

\begin{center}
\footnotesize
\begin{tabular}{lrrrr}
\toprule
 & \multicolumn{2}{c}{Corta ($N=248$)} & \multicolumn{2}{c}{Larga ($N=293$)} \\
Variable & Media & Desv. & Media & Desv. \\
\midrule
EMBI (pb)     & 213,12 & 78,41 & 221,49 & 86,16 \\
JLoss         & 6,05   & 8,79  & 5,87   & 8,26  \\
GaR (q05)     & 0,003  & 0,041 & 0,003  & 0,039 \\
Deuda/PIB (\%)& 52,39  & 21,78 & 45,42  & 24,09 \\
Bal. fiscal/PIB (\%) & $-3,21$ & 2,66 & $-3,20$ & 2,65 \\
Reservas/PIB (\%)    & 18,62 & 7,04 & 18,56 & 7,05 \\
C. corriente/PIB (\%)& $-2,90$ & 1,71 & $-2,86$ & 1,72 \\
Inflación YoY (\%)   & 4,13 & 1,81 & 4,12 & 1,81 \\
REER (índice)        & 116,22 & 19,51 & 116,39 & 19,62 \\
\bottomrule
\end{tabular}

\medskip
\footnotesize\textbf{Cuadro 2.} Estadísticos descriptivos, núcleo LatAm corta vs. larga.
\end{center}

La heterogeneidad transversal documentada en la base corta (Brasil concentra la fragilidad bancaria del panel, con JLoss medio de 17,5 frente a 2,4–4,7 en el resto) se mantiene esencialmente inalterada en la base larga (JLoss medio de Brasil: 16,1), y la descomposición de varianza \textit{within/total} —crítica para un estimador de efectos fijos— es prácticamente idéntica entre ambas bases: el GaR conserva un cociente \textit{within/total} de 0,95 y el JLoss de 0,75–0,78, mientras que la deuda/PIB y las reservas/PIB siguen siendo mayoritariamente transversales (0,37–0,41). La fuente de identificación del modelo no cambia al extender la ventana.

\subsection{Cobertura efectiva y la extensión a 11 países}

La cobertura nominal de cada base no coincide exactamente con la cobertura \textbf{efectiva} de las regresiones: el GaR viene indefinido en los primeros trimestres de entrenamiento de cada país (tamaño de muestra insuficiente para la regresión cuantílica) y esas filas se excluyen automáticamente vía \texttt{na.omit}. En la base larga, Colombia y Perú pierden en efecto sus primeros trimestres (55 de 59 y 51 de 59 observaciones válidas, respectivamente); en la base extendida, Bulgaria pierde la mitad de sus 16 trimestres nominales, quedando con apenas 8 observaciones utilizables —tratable como anecdótico, no como evidencia con poder estadístico.

La Figura 1 grafica la cobertura temporal de las tres bases por país. Se aprecia que la base extendida agrega países casi exclusivamente en la ventana 2010Q1–2014Q4 (con la excepción de los cinco países LatAm, presentes en toda la ventana): Sudáfrica aporta solo 6 trimestres, Bulgaria 16 (8 efectivos), y Bulgaria/China/Indonesia/Polonia/Turquía 20 cada uno. Esta estructura —un núcleo con $T$ largo rodeado de países con $T$ corto— se explota en la Sección 6.12 restringiendo el análisis a la ventana común, la comparación más simétrica disponible entre el núcleo LatAm y la extensión.

\begin{center}
\includegraphics[width=\linewidth]{fig_cobertura.pdf}
\footnotesize\textbf{Figura 1.} Cobertura temporal por base y país.
\end{center}

En la base extendida (Cuadro 3), el EMBI medio de los seis países no-LatAm (216,5 pb en el agregado de 11 países) es comparable en magnitud al del núcleo LatAm, pero la fragilidad bancaria media es sustancialmente menor y más homogénea (JLoss entre 2,3 y 5,1 en los seis países nuevos, frente a la asimetría dominada por Brasil en el núcleo), un contraste relevante para interpretar el hallazgo exploratorio de la Sección 6.12.

\begin{center}
\footnotesize
\begin{tabular}{lrrr}
\toprule
País & EMBI (pb) & JLoss & $T$ \\
\midrule
Bulgaria     & 195,0 & 2,96 & 16 \\
China        & 164,6 & 2,32 & 20 \\
Indonesia    & 243,2 & 2,84 & 20 \\
Polonia      & 161,5 & 2,81 & 20 \\
Sudáfrica    & 165,3 & 2,68 & 6  \\
Turquía      & 256,3 & 5,09 & 20 \\
\bottomrule
\end{tabular}

\medskip
\footnotesize\textbf{Cuadro 3.} Medias por país, seis economías no-LatAm de la base extendida.
\end{center}

% ---

\section{Resultados}

\subsection{Estrategia de identificación y lectura de la evidencia}

La estimación se organiza en torno al núcleo LatAm descrito en la Sección 5 ($N=248$ en su ventana original; $T$ entre 48 y 50). La especificación de referencia incorpora efectos fijos bidireccionales —de país y de tiempo—, de modo que los efectos de país absorben la heterogeneidad estructural no observada y los efectos de tiempo capturan la totalidad de los choques globales comunes. La inferencia se basa en errores estándar de Driscoll–Kraay. La hipótesis central —complementariedad entre fragilidad bancaria sistémica y riesgo de cola izquierda del crecimiento— se contrasta a través del término de interacción $\theta$ en

\[EMBI_{i,t}=\alpha_i+\lambda_t+\beta_1 JLoss_{i,t}+\beta_2 GaR_{i,t}+\theta (JLoss\times GaR)_{i,t}+\gamma' X_{i,t}+\varepsilon_{i,t}.\]

Bajo la convención de signos adoptada, la predicción teórica del mecanismo de \textit{doom loop} es inequívoca: $\theta<0$.

\subsection{Resultado principal: el canal de complementariedad}

El coeficiente de interacción $\hat\theta=-0{,}352$ (M2, sin controles) es negativo y altamente significativo ($t=-3{,}97$), precisamente el signo predicho por el multiplicador de \textit{doom loop}. La introducción de la interacción reduce el coeficiente lineal de JLoss (de 2,73 a 1,23): el efecto promedio se descompone en un componente de nivel y uno dependiente del estado del riesgo de cola.

M3 añade el vector completo de controles macroeconómicos domésticos. El resultado crucial es la estabilidad cuantitativa de la interacción: $\hat\theta$ pasa de $-0{,}352$ a $-0{,}316$ y conserva significancia al 1\% ($t=-3{,}09$). Que la magnitud del parámetro de interés permanezca prácticamente inalterada tras condicionar por deuda, balance fiscal, reservas, cuenta corriente, inflación y tipo de cambio real es evidencia de que la complementariedad no es un artefacto de variables omitidas de naturaleza fiscal o externa.

\subsection{Triangulación: ventana temporal extendida (corta vs. larga)}

Con el GaR reentrenado sobre el panel de 15 países —que aporta más historia de entrenamiento y extiende la ventana estimable hasta 2007Q4—, se reestima la batería completa sobre la base larga y se compara lado a lado con la corta (Cuadro 4).

\begin{table*}[t]
\centering
\small
\begin{tabular}{lrrrrr}
\toprule
 & \multicolumn{1}{c}{M1} & \multicolumn{1}{c}{M2} & \multicolumn{1}{c}{M3} & \multicolumn{1}{c}{M4 (log)} & \multicolumn{1}{c}{M5 (FE país)} \\
\midrule
\multicolumn{6}{l}{\textit{Panel A: LatAm corta ($N=248/200$)}} \\
JLoss        & 2,731*** & 1,233*** & 1,340* & 0,007*** & 0,171 \\
GaR (pp)     & 3,423*   & 2,863    & 1,663  & 0,009    & $-4,408^{**}$ \\
JLoss$\times$GaR & — & $-0,352^{***}$ & $-0,316^{***}$ & $-0,000$ & $-0,223^{*}$ \\
\midrule
\multicolumn{6}{l}{\textit{Panel B: LatAm larga ($N=293/253$)}} \\
JLoss        & 2,261*** & 1,271*** & 1,580*** & 0,006*** & 0,605 \\
GaR (pp)     & 1,712    & 1,558    & 0,893    & 0,004    & $-5,056^{**}$ \\
JLoss$\times$GaR & — & $-0,336^{***}$ & $-0,305^{**}$ & $-0,000$ & $-0,226^{*}$ \\
\bottomrule
\end{tabular}
\smallskip
\small\textbf{Cuadro 4.} Núcleo LatAm — corta (2010Q1–2022Q2) vs. larga (2007Q4–2022Q2). SE de Driscoll–Kraay entre paréntesis omitidos por espacio; *** $p<0{,}01$, ** $p<0{,}05$, * $p<0{,}10$. M1–M4: FE país+tiempo. M5: FE país + VIX explícito.
\end{table*}

El resultado es inequívoco: extender el GaR con más historia de entrenamiento —el ejercicio de robustez más directo disponible sobre la propia construcción de la métrica— no altera la conclusión central. $\hat\theta$ se mantiene negativo y significativo al 1--5\% en M2--M5 de ambas bases, con magnitudes prácticamente indistinguibles ($-0{,}352$ a $-0{,}316$ en la corta; $-0{,}336$ a $-0{,}226$ en la larga). La Figura 2 superpone el efecto marginal $\partial EMBI/\partial JLoss$ evaluado sobre el soporte de GaR para ambas bases: las dos curvas son estadísticamente indistinguibles, con la banda de confianza al 95\% cruzando cero en la misma región del riesgo de cola.

\begin{center}
\includegraphics[width=\linewidth]{fig_efecto_marginal.pdf}
\footnotesize\textbf{Figura 2.} Efecto marginal de JLoss sobre el spread según severidad de GaR, corta vs. larga (M3 de cada base).
\end{center}

Los diagnósticos son igualmente consistentes entre bases: la prueba de Pesaran rechaza la independencia transversal en ambas ($z=-4{,}95$ corta, $z=-5{,}20$ larga; $p<0{,}001$), la prueba de Wooldridge detecta autocorrelación serial en ambas ($\chi^2=198{,}9$ y $207{,}2$; $p<0{,}001$), y el test de Hausman rechaza la consistencia de efectos aleatorios en ambas ($\chi^2=16{,}0$ y $15{,}9$; $p<0{,}01$) —validando la misma estrategia de inferencia (Driscoll--Kraay) y de efectos fijos en las dos ventanas. Los factores de inflación de varianza permanecen uniformemente bajos (máximo 3,33 en la corta, 2,37 en la larga).

\subsection{Robustez con los controles del paper de referencia como índice agregado}

Los seis controles adicionales que replican los del paper de referencia de JLoss (Chari et al., 2024) —deuda/PIB (log), margen de utilidad bancaria, tasa del Tesoro a 10 años (log), \textit{spread} de alto rendimiento (log) y \textit{spread} on/off-the-run (log)— se incorporan como \textbf{dos índices agregados} en vez de regresores sueltos, por dos razones: primero, el margen de utilidad bancaria solo tiene cobertura real reconstruida para 1999–2011; segundo, y más importante, mezclar variables de variación transversal-temporal (deuda/PIB, margen de utilidad) con variables de variación puramente temporal (tasa del Tesoro, \textit{spread} on/off-the-run —idénticas para todos los países en cada trimestre) bajo efectos fijos bidireccionales es econométricamente incómodo: la parte global de la mezcla queda casi enteramente absorbida por los efectos de tiempo, y el índice combinado pierde poder de forma artificial (se verificó empíricamente: con un índice único y FE país+tiempo, el $R^2$ \textit{within} se torna incluso negativo en la base corta). La solución —la misma lógica que ya emplea M5 en el Cuadro 4— es separar:

\begin{itemize}
    \item $X^{dom}_{i,t}$ = promedio de \textit{z-scores} de deuda/PIB (log) y margen de utilidad bancaria, usado con \textbf{efectos fijos país+tiempo}.
    \item $X^{global}_{i,t}$ = promedio de \textit{z-scores} de la tasa del Tesoro a 10 años (log) y el \textit{spread} on/off-the-run (log), usado con \textbf{efectos fijos de país únicamente}, replicando el argumento de identificación de M5.
\end{itemize}

El \textit{spread} de alto rendimiento queda excluido de ambos índices: el proxy descargado (ICE BofA US HY OAS) llegó sin cobertura útil en la ventana de estudio. La volatilidad cambiaria del paper de referencia permanece pendiente de construcción.

\begin{table*}[t]
\centering
\small
\begin{tabular}{lrrrrr}
\toprule
 & \multicolumn{2}{c}{$X^{dom}$ (FE2)} & \multicolumn{3}{c}{$X^{global}$ (FE país)} \\
 & Corta & Larga & Corta & Larga & Ext. \\
\midrule
JLoss & 1,866*** & 1,455*** & 1,066 & 1,643** & 2,405*** \\
GaR (pp) & 3,494 & 1,180 & $-7,067^{***}$ & $-5,118^{***}$ & $-5,141^{***}$ \\
JLoss$\times$GaR & $-0,306^{***}$ & $-0,327^{***}$ & $-0,358^{***}$ & $-0,391^{**}$ & $-0,157$ \\
$X^{dom}$ & $-46,674$ & $-16,274$ & — & — & — \\
$X^{global}$ & — & — & $-13,949$ & $-62,191^{***}$ & $-57,288^{***}$ \\
$N$ & 206 & 280 & 248 & 281 & 374 \\
\bottomrule
\end{tabular}
\smallskip
\small\textbf{Cuadro 5.} $\theta$ con los controles del paper de referencia como índice agregado. SE de Driscoll–Kraay; *** $p<0{,}01$, ** $p<0{,}05$, * $p<0{,}10$.
\end{table*}

$\theta$ se mantiene negativo y significativo en cuatro de las cinco columnas —el resultado central sobrevive a agregar los controles del paper de referencia, tanto en su versión doméstica como global, en la base corta y en la larga—. El índice global resulta además significativo por sí mismo bajo efectos fijos de país (a diferencia de la versión mezclada bajo FE país+tiempo), confirmando que la separación doméstico/global fue la decisión correcta y no una mera formalidad. En la base extendida, $\theta$ conserva el signo pero pierde significancia bajo $X^{global}$ ($p=0{,}37$) —anticipando el patrón de la Sección 6.12— y $X^{dom}$ se omite de la tabla porque su $N$ cae a 83 (solo los tramos LatAm con margen de utilidad real, 2007–2011), sin interpretación a nivel de las 11 economías.

\subsection{Generalización fuera de LatAm: panel extendido a 11 países}

La prueba más directa de si la complementariedad JLoss$\times$GaR se generaliza fuera del núcleo latinoamericano consiste en restringir el panel extendido a la ventana en la que efectivamente coexisten múltiples países (2010Q1–2014Q4, 11 economías) y comparar el coeficiente de interacción con el que arroja el subpanel de solo-LatAm en esa \textbf{misma} ventana —la comparación de manzanas con manzanas más limpia disponible con los datos actuales.

\begin{center}
\footnotesize
\begin{tabular}{lrrr}
\toprule
 & \shortstack{Ext.\\FE país} & \shortstack{Ext.\\FE2} & \shortstack{Solo LatAm-5\\FE2} \\
\midrule
JLoss & 3,386*** & 2,376*** & 1,464*** \\
GaR (pp) & $-4,697^{**}$ & $-3,587$ & 0,365 \\
JLoss$\times$GaR & 1,072** & 0,626** & $-0,074$ \\
$N$ & 193 & 193 & 100 \\
\bottomrule
\end{tabular}

\medskip
\textbf{Cuadro 6.} Ventana común 2010Q1–2014Q4: extendido (11 países) vs. solo LatAm-5. SE de Driscoll–Kraay.
\end{center}

El resultado debe reportarse con honestidad, no maquillarse: en esta ventana, el subpanel de solo-LatAm no rechaza $\theta=0$ ($\hat\theta=-0{,}074$, $p=0{,}56$), mientras que el panel de 11 países en la misma ventana entrega un coeficiente \textbf{positivo} y marginalmente significativo ($\hat\theta=+0{,}626$ a $+1{,}072$, $p\approx0{,}02$–$0{,}05$) —signo \textbf{contrario} a la hipótesis de complementariedad. Esto no invalida el resultado del núcleo LatAm, que descansa sobre $T$ largo y controles domésticos completos (Secciones 6.2–6.4); pero sí establece que la complementariedad \textbf{no se generaliza mecánicamente} fuera de América Latina con los datos disponibles hoy. La muestra extendida combina países con $T$ muy corto (Sudáfrica $n\approx6$, Bulgaria $n\approx8$ efectivo), dos fuentes de EMBI distintas (tablas del FMI vs. BCRP) y estructuras económicas heterogéneas —China, Turquía, Polonia, Indonesia, Sudáfrica y Bulgaria en un solo \textit{pool}— sin controles domésticos que absorban esas diferencias estructurales. El signo positivo es, por tanto, más creíble como síntoma de heterogeneidad no modelada que como evidencia de que la relación se invierte fuera de LatAm; con $T$ entre 6 y 20 para los países no-LatAm, tanto la prueba de Pesaran como Driscoll–Kraay operan además fuera de su rango ideal de aplicación. La recomendación metodológica que se sigue en este trabajo es anclar el resultado principal en el núcleo LatAm y presentar esta sección como evidencia exploratoria que motiva trabajo futuro —más $T$ por país, controles domésticos para las seis economías nuevas— antes de generalizar la conclusión.

\subsection{Síntesis: triangulación de $\theta$ a través de bases y especificaciones}

La Figura 3 sintetiza el coeficiente $\theta$ y su intervalo de confianza al 95\% a través de las doce especificaciones estimadas en las Secciones 6.2–6.5.

\begin{center}
\includegraphics[width=\linewidth]{fig_forest_theta.pdf}
\footnotesize\textbf{Figura 3.} Triangulación del coeficiente de interacción entre bases y especificaciones.
\end{center}

La lectura de conjunto es clara. Dentro del núcleo LatAm —ocho especificaciones entre la base corta y la larga, con y sin los controles del paper de referencia— $\theta$ es sistemáticamente negativo, con magnitudes entre $-0{,}31$ y $-0{,}39$ y significancia estadística en todos los casos. Fuera de ese núcleo, la evidencia es mixta: el panel extendido en su muestra completa (desbalanceada) mantiene el signo negativo pero pierde significancia ($p\approx0{,}32$–$0{,}37$), mientras que la ventana común 2010–2014 exhibe el único signo positivo del conjunto, y precisamente en el subpanel más heterogéneo. Esta estructura de resultados —robustez dentro de LatAm, fragilidad fuera de ella— es en sí misma informativa: delimita con precisión lo que la evidencia actual permite afirmar y lo que exige más datos antes de poder afirmarse.

\subsection{Magnitud económica: el multiplicador del riesgo soberano}

El gradiente $\partial EMBI/\partial JLoss = \hat\beta_1+\hat\theta(GaR-\overline{GaR})$, evaluado sobre M3 del núcleo LatAm, es monótono y económicamente sustancial (Cuadro 7): cuando el riesgo de cola es severo (percentil 10), una unidad adicional de fragilidad bancaria eleva el spread en 2,5 pb ($t=4{,}59$); en el centro de la distribución el impacto se reduce a la mitad (1,1 pb); en entornos de crecimiento benigno (percentil 90) se torna estadísticamente nulo. Un aumento de una desviación estándar intra-país de JLoss ($\approx6{,}6$ unidades) se traduce en un encarecimiento de unos 17 pb del financiamiento soberano bajo estrés de cola, frente a un efecto indistinguible de cero en condiciones favorables. Este patrón es la firma empírica de la amplificación supra-aditiva: la fragilidad bancaria y el riesgo de cola del crecimiento no se suman, se potencian.

\begin{center}
\footnotesize
\begin{tabular}{lrrr}
\toprule
Estado del GaR & Nivel (pp) & $\partial EMBI/\partial JLoss$ (pb) & $t$ \\
\midrule
Cola severa (p10)   & $-3,57$ & $+2,509$ & 4,59 \\
Mediana              & $+0,37$ & $+1,078$ & 1,85 \\
Entorno benigno (p90)& $+4,95$ & $-0,589$ & $-0,65$ \\
\bottomrule
\end{tabular}

\medskip
\footnotesize\textbf{Cuadro 7.} Efecto marginal de JLoss según severidad del riesgo de cola (núcleo LatAm corta, M3).
\end{center}

\subsection{No linealidad: el modelo de umbral de Hansen}

Como contraste no lineal alternativo, se estima un modelo de umbral de panel (Hansen, 1999, 2000) sobre el núcleo LatAm. El umbral estimado es $\hat\gamma=2{,}76$ pp (IC95\% $[-3{,}57,\,4{,}95]$). En el régimen de cola severa ($GaR\le\hat\gamma$), el efecto de la fragilidad sobre el spread es de $+2{,}76$ pb por unidad; en el régimen benigno ($GaR>\hat\gamma$) es de $-1{,}19$ pb, indistinguible de cero. El modelo de umbral, que no impone la forma funcional de la interacción, corrobora así el hallazgo del término multiplicativo.

\subsection{Robustez adicional del núcleo LatAm}

La significancia de $\theta$ es invariante al estimador de errores estándar (Driscoll–Kraay, \textit{clustering} por país, por tiempo, HAC de Arellano: $t$ entre $-3{,}41$ y $-3{,}91$ en la corta). Excluyendo la pandemia, la interacción se intensifica en ambas bases ($\hat\theta=-0{,}255$, $t=-3{,}17$ en la corta sin COVID; $\hat\theta=-0{,}383$, $t=-3{,}37$ en la larga sin COVID). En una especificación dinámica con el spread rezagado, la persistencia es alta (coeficiente autorregresivo $\approx0{,}85$ en ambas bases) y la interacción se atenúa pero conserva signo correcto ($\hat\theta=-0{,}120$, $p=0{,}09$ en la corta; $\hat\theta=-0{,}105$, $p=0{,}15$ en la larga) —diagnóstico de persistencia, no estimación preferida, dado el sesgo de Nickell bajo efectos fijos con $T$ moderado—. Sustituyendo el cuantil 5\% por el \textit{Expected Shortfall}, la interacción conserva signo negativo y significancia al 1\% ($\hat\theta=-0{,}262$, $t=-2{,}66$). Reestimando $\theta$ excluyendo secuencialmente cada país, el coeficiente permanece negativo y significativo en las cinco submuestras.

\subsection{Los controles macroeconómicos y diagnósticos del panel}

El tipo de cambio real efectivo se comporta de manera ejemplar: negativo y altamente significativo en todas las especificaciones y en ambas bases (corta M3: $-1{,}05$, $t=-3{,}67$; larga M3: $-1{,}09$, $t=-4{,}17$). Tres controles fiscales y externos exhiben signos contraintuitivos (balance fiscal y cuenta corriente positivos, deuda/PIB débilmente negativo) en ambas bases, consistentemente con endogeneidad por causalidad inversa —no con multicolinealidad, dado que los VIF son uniformemente bajos—. El test $F$ de efectos de tiempo establece su significancia conjunta en ambas bases, validando la estructura bidireccional.

\subsection{Síntesis}

La evidencia respalda de manera consistente la hipótesis de complementariedad \textbf{dentro del núcleo latinoamericano}: el coeficiente de interacción es negativo, significativo y cuantitativamente estable a través de dos ventanas temporales que difieren en casi tres años de historia (2010Q1 vs. 2007Q4), de la inclusión progresiva de controles domésticos, de los controles del paper de referencia agregados como índice, de cuatro estimadores de la matriz de varianzas, de la exclusión del episodio pandémico, de dos medidas alternativas del riesgo de cola, de la exclusión secuencial de cada país y de un modelo de umbral que no impone la forma funcional de la interacción. La extensión exploratoria a 11 economías emergentes, en cambio, no reproduce el mismo patrón de forma limpia: el signo se mantiene en la muestra completa pero pierde significancia, y se invierte —de forma marginalmente significativa— en la ventana común más comparable. Esta asimetría entre robustez interna y fragilidad externa es, en sí misma, el resultado más honesto que la evidencia actual permite reportar: lo que determina la prima de riesgo soberano en el núcleo de economías latinoamericanas estudiado no es la fragilidad bancaria ni el riesgo de cola por separado, sino su interacción; si ese mecanismo opera de la misma manera en el resto de las economías emergentes es, con los datos hoy disponibles, una pregunta abierta y no una conclusión establecida.

% ---

\section{Discusión y Conclusiones}

\subsection{Interpretación económica}

El hallazgo central de este trabajo —un coeficiente de interacción JLoss$\times$GaR negativo, significativo y estable dentro del núcleo LatAm, replicado en dos ventanas temporales y robusto a los controles del paper de referencia— admite una interpretación económica directa en términos del mecanismo de \textit{doom loop}. La fragilidad del sistema bancario opera como una aproximación al pasivo contingente que el soberano podría verse forzado a asumir mediante un rescate; pero ese pasivo sólo se torna relevante para la prima de riesgo cuando la economía transita un estado de cola. La evidencia de que el efecto de la fragilidad bancaria pasa de ser nulo en entornos benignos a $+2{,}5$ pb por unidad bajo estrés severo es la firma empírica del multiplicador: no la suma de dos vulnerabilidades, sino su producto.

El hallazgo exploratorio de la Sección 6.12 —la no replicación limpia fuera de LatAm— merece una lectura igualmente directa: no como una anomalía a ocultar, sino como evidencia de que el mecanismo, tal como está identificado hoy, requiere condiciones de datos (T largo, controles domésticos completos, fuente de EMBI homogénea) que actualmente solo se satisfacen para el núcleo latinoamericano.

\subsection{Implicancias de política}

Si la transmisión de la fragilidad bancaria al costo de financiamiento soberano se concentra en los estados de cola del crecimiento, la regulación orientada a la resiliencia del sistema bancario —capital, liquidez, provisiones contracíclicas— reduce la exposición del soberano a los escenarios macroeconómicos adversos y su prima de riesgo en los mercados internacionales. La fragilidad bancaria y el riesgo de cola del crecimiento deben monitorearse conjuntamente: la coincidencia de ambos, y no cada uno por separado, es la señal de alarma relevante. Esta implicancia está bien establecida para economías latinoamericanas de tamaño y desarrollo financiero medio; su extensión a otras economías emergentes es, a la luz de la Sección 6.12, una hipótesis de trabajo y no todavía una recomendación generalizable.

\subsection{Limitaciones}

Primero, el resultado robusto de esta entrega descansa en cinco economías; la extensión exploratoria a 11 países no lo contradice pero tampoco lo confirma, y la ampliación al panel completo de $\sim$19 economías —condición para una validez externa genuina— sigue pendiente de completar la extracción de JLoss (Hungría, India, Corea del Sur sin trabajo iniciado; Turquía, Sudáfrica, Polonia, Bulgaria, Indonesia, Pakistán y China con extractor listo pero no ejecutado en un entorno con acceso a red) y de EMBI real más allá de 2014 para los seis países no-LatAm (recuperado hoy solo de tablas del FMI 2011–2015). Segundo, los controles fiscales y externos exhiben signos compatibles con endogeneidad por causalidad inversa, no corregida en este diseño. Tercero, la elevada persistencia del spread sugiere una dimensión dinámica no modelada por la especificación estática, sujeta al sesgo de Nickell bajo efectos fijos. Cuarto, dos controles del paper de referencia permanecen incompletos: el margen de utilidad bancaria solo está reconstruido hasta 2011, y la volatilidad cambiaria no se ha construido todavía; el \textit{spread} de alto rendimiento (proxy ICE BofA) se descargó sin cobertura útil en la ventana de estudio y se excluyó de los índices agregados. Quinto, las métricas centrales descansan en supuestos estructurales —la aproximación de punto de silla bajo independencia condicional a un único factor sistémico, en el caso de JLoss; la especificación de la regresión cuantílica, en el de GaR— cuya robustez, si bien validada internamente, condiciona la inferencia.

\subsection{Agenda futura}

De las limitaciones anteriores se desprende una agenda concreta, ya en curso: (i) ejecutar en un entorno con acceso a red los siete extractores de JLoss ya escritos (Turquía, Sudáfrica —con los XML de balance ya descargados—, Polonia, Bulgaria, Indonesia, Pakistán, China) y desarrollar los tres pendientes (Hungría vía MNB, India vía RBI, Corea del Sur vía FSS/DART, fuentes ya identificadas); (ii) buscar ediciones del \textit{Global Financial Stability Report} posteriores a 2015 para extender el EMBI real de los países no-LatAm más allá de 2014, o evaluar acceso institucional a Bloomberg/Refinitiv/CEIC; (iii) resolver la descarga del \textit{spread} de alto rendimiento y construir la volatilidad cambiaria para completar los seis controles del paper de referencia; (iv) con esos datos, repetir la comparación de la Sección 6.12 con un panel balanceado y controles domésticos completos para las seis economías nuevas —la prueba de generalización genuinamente informativa que los datos actuales no permiten realizar todavía—; y (v) el tratamiento formal de la dinámica y la endogeneidad mediante estimadores GMM dinámicos (Arellano y Bond, 1991).

\subsection{Conclusión}

Este trabajo articula, en un único marco empírico, dos literaturas que la investigación previa había mantenido separadas: la del nexo banca-soberano y la del riesgo a la baja del crecimiento. La evidencia, triangulada sobre tres bases de datos y dos ventanas temporales, respalda de manera consistente y robusta la hipótesis central \textbf{dentro del núcleo de cinco economías latinoamericanas}: la fragilidad bancaria sistémica y el riesgo de cola del crecimiento no determinan el spread soberano de forma aditiva e independiente, sino complementaria. Esa misma triangulación, aplicada con honestidad a la extensión exploratoria fuera de LatAm, muestra que la generalización del mecanismo es todavía una pregunta abierta, no una conclusión establecida —una delimitación que no debilita el resultado central, sino que precisa exactamente su alcance actual y traza la agenda de datos necesaria para extenderlo. Todos los archivos computacionales ocupados para la reproducción de esta investigación se encuentran en el repositorio del proyecto.

\section*{Referencias (selección preliminar)}

\begin{itemize}
    \item Acharya, V., Drechsler, I., \& Schnabl, P. (2014). A Pyrrhic Victory? Bank Bailouts and Sovereign Credit Risk. \textit{Journal of Finance}, 69(6), 2689–2739.
    \item Acharya, V. V., Pedersen, L. H., Philippon, T., \& Richardson, M. (2017). Measuring Systemic Risk. \textit{Review of Financial Studies}, 30(1), 2–47.
    \item Acharya, V. V., \& Steffen, S. (2015). The ``Greatest'' Carry Trade Ever? Understanding Eurozone Bank Risks. \textit{Journal of Financial Economics}, 115(2), 215–236.
    \item Adrian, T., \& Brunnermeier, M. K. (2016). CoVaR. \textit{American Economic Review}, 106(7), 1705–1741.
    \item Adrian, T., Boyarchenko, N., \& Giannone, D. (2019). Vulnerable Growth. \textit{American Economic Review}, 109(4), 1263–1289.
    \item Allen, L., Bali, T. G., \& Tang, Y. (2012). Does Systemic Risk in the Financial Sector Predict Future Economic Downturns? \textit{Review of Financial Studies}, 25(10), 3000–3036.
    \item Arellano, M., \& Bond, S. (1991). Some Tests of Specification for Panel Data: Monte Carlo Evidence and an Application to Employment Equations. \textit{Review of Economic Studies}, 58(2), 277–297.
    \item Barndorff-Nielsen, O. E., \& Shephard, N. (2002). Econometric Analysis of Realized Volatility and Its Use in Estimating Stochastic Volatility Models. \textit{Journal of the Royal Statistical Society B}, 64(2), 253–280.
    \item Bank for International Settlements (BIS). (2006). \textit{International Convergence of Capital Measurement and Capital Standards (Basel II)}. Basel.
    \item Billio, M., Getmansky, M., Lo, A. W., \& Pelizzon, L. (2012). Econometric Measures of Connectedness and Systemic Risk in the Finance and Insurance Sectors. \textit{Journal of Financial Economics}, 104(3), 535–559.
    \item Bisias, D., Flood, M., Lo, A. W., \& Valavanis, S. (2012). A Survey of Systemic Risk Analytics. \textit{Annual Review of Financial Economics}, 4, 255–296.
    \item Bocola, L. (2016). The Pass-Through of Sovereign Risk. \textit{Journal of Political Economy}, 124(4), 879–926.
    \item Bolton, P., \& Jeanne, O. (2011). Sovereign Default Risk and Bank Fragility in Financially Integrated Economies. \textit{IMF Economic Review}, 59(2), 162–194.
    \item Brownlees, C., \& Engle, R. (2016). SRISK: A Conditional Capital Shortfall Measure of Systemic Risk. \textit{Review of Financial Studies}, 30(1), 48–79.
    \item Brunnermeier, M. K., Langfield, S., Pagano, M., Reis, R., Van Nieuwerburgh, S., \& Vayanos, D. (2016). ESBies: Safety in the Tranches. \textit{Economic Policy}, 32(90), 175–219.
    \item Brunnermeier, M. K., \& Sannikov, Y. (2014). A Macroeconomic Model with a Financial Sector. \textit{American Economic Review}, 104(2), 379–421.
    \item Calvo, G. A., Leiderman, L., \& Reinhart, C. M. (1996). Inflows of Capital to Developing Countries in the 1990s. \textit{Journal of Economic Perspectives}, 10(2), 123–139.
    \item Canay, I. A. (2011). A Simple Approach to Quantile Regression for Panel Data. \textit{Econometrics Journal}, 14(3), 368–386.
    \item Cavallo, E. A., \& Valenzuela, P. (2010). The Determinants of Corporate Risk in Emerging Markets: An Option-Adjusted Spread Analysis. \textit{International Journal of Finance \& Economics}, 15(1), 59–74.
    \item Chari, A., Garcés, F., Martínez, J. F., \& Valenzuela, P. (2024). Sovereign Credit Spreads, Banking Fragility, and Global Factors. \textit{Journal of Financial Stability}, 72, 101235.
    \item Chernozhukov, V., Fernández-Val, I., \& Galichon, A. (2010). Quantile and Probability Curves Without Crossing. \textit{Econometrica}, 78(3), 1093–1125.
    \item Choi, I. (2001). Unit Root Tests for Panel Data. \textit{Journal of International Money and Finance}, 20(2), 249–272.
    \item Csonto, B., \& Ivaschenko, I. (2013). Determinants of Sovereign Bond Spreads in Emerging Markets. \textit{IMF Working Paper} 13/164.
    \item Dieckmann, S., \& Plank, T. (2012). Default Risk of Advanced Economies: An Empirical Analysis of CDS during the Financial Crisis. \textit{Review of Finance}, 16(4), 903–934.
    \item Driscoll, J. C., \& Kraay, A. C. (1998). Consistent Covariance Matrix Estimation with Spatially Dependent Panel Data. \textit{Review of Economics and Statistics}, 80(4), 549–560.
    \item Edwards, S. (1984). LDC Foreign Borrowing and Default Risk: An Empirical Investigation, 1976–80. \textit{American Economic Review}, 74(4), 726–734.
    \item Farhi, E., \& Tirole, J. (2018). Deadly Embrace: Sovereign and Financial Balance Sheets Doom Loops. \textit{Review of Economic Studies}, 85(3), 1781–1823.
    \item Fratzscher, M., \& Rieth, M. (2019). Monetary Policy, Bank Bailouts and the Sovereign-Bank Risk Nexus in the Euro Area. \textit{Review of Finance}, 23(4), 745–775.
    \item Gennaioli, N., Martin, A., \& Rossi, S. (2014). Sovereign Default, Domestic Banks, and Financial Institutions. \textit{Journal of Finance}, 69(2), 819–866.
    \item Giglio, S., Kelly, B., \& Pruitt, S. (2016). Systemic Risk and the Macroeconomy: An Empirical Evaluation. \textit{Journal of Financial Economics}, 119(3), 457–471.
    \item González-Rozada, M., \& Levy-Yeyati, E. (2008). Global Factors and Emerging Market Spreads. \textit{Economic Journal}, 118(533), 1917–1936.
    \item Hansen, B. E. (1999). Threshold Effects in Non-Dynamic Panels: Estimation, Testing, and Inference. \textit{Journal of Econometrics}, 93(2), 345–368.
    \item Hansen, B. E. (2000). Sample Splitting and Threshold Estimation. \textit{Econometrica}, 68(3), 575–603.
    \item Hatzius, J., Hooper, P., Mishkin, F. S., Schoenholtz, K. L., \& Watson, M. W. (2010). Financial Conditions Indexes: A Fresh Look after the Financial Crisis. \textit{NBER Working Paper} 16150.
    \item He, Z., \& Krishnamurthy, A. (2013). Intermediary Asset Pricing. \textit{American Economic Review}, 103(2), 732–770.
    \item Hilscher, J., \& Nosbusch, Y. (2010). Determinants of Sovereign Risk: Macroeconomic Fundamentals and the Pricing of Sovereign Debt. \textit{Review of Finance}, 14(2), 235–262.
    \item Holló, D., Kremer, M., \& Lo Duca, M. (2012). CISS – A Composite Indicator of Systemic Stress in the Financial System. \textit{ECB Working Paper} 1426.
    \item Huang, X., Zhou, H., \& Zhu, H. (2009). A Framework for Assessing the Systemic Risk of Major Financial Institutions. \textit{Journal of Banking \& Finance}, 33(11), 2036–2049.
    \item Imbens, G. W., \& Wooldridge, J. M. (2009). Recent Developments in the Econometrics of Program Evaluation. \textit{Journal of Economic Literature}, 47(1), 5–86.
    \item Kealhofer, S. (2000). The Quantification of Credit Risk. \textit{KMV Corporation} (mimeo).
    \item Koenker, R. (2004). Quantile Regression for Longitudinal Data. \textit{Journal of Multivariate Analysis}, 91(1), 74–89.
    \item Koenker, R. (2005). \textit{Quantile Regression}. Econometric Society Monographs. Cambridge University Press.
    \item Koenker, R., \& Bassett, G. (1978). Regression Quantiles. \textit{Econometrica}, 46(1), 33–50.
    \item Laeven, L., \& Valencia, F. (2013). Systemic Banking Crises Database. \textit{IMF Economic Review}, 61(2), 225–270.
    \item Laeven, L., \& Valencia, F. (2020). Systemic Banking Crises Database II. \textit{IMF Economic Review}, 68(2), 307–361.
    \item Leonello, A. (2018). Government Guarantees and the Two-Way Feedback between Banking and Sovereign Debt Crises. \textit{Journal of Financial Economics}, 130(3), 592–619.
    \item Longstaff, F. A., Pan, J., Pedersen, L. H., \& Singleton, K. J. (2011). How Sovereign Is Sovereign Credit Risk? \textit{American Economic Journal: Macroeconomics}, 3(2), 75–103.
    \item Machado, J. A. F., \& Santos Silva, J. M. C. (2019). Quantiles via Moments. \textit{Journal of Econometrics}, 213(1), 145–173.
    \item Martin, R., Thompson, K., \& Browne, C. (2001). Taking to the Saddle. \textit{Risk}, 14(6), 91–94.
    \item Merton, R. C. (1974). On the Pricing of Corporate Debt. \textit{Journal of Finance}, 29(2), 449–470.
    \item Miranda-Agrippino, S., \& Rey, H. (2022). The Global Financial Cycle. \textit{Handbook of International Economics}, Vol. 6.
    \item Mody, A., \& Sandri, D. (2012). The Eurozone Crisis: How Banks and Sovereigns Came to Be Joined at the Hip. \textit{Economic Policy}, 27(70), 199–230.
    \item Nickell, S. (1981). Biases in Dynamic Models with Fixed Effects. \textit{Econometrica}, 49(6), 1417–1426.
    \item Ossandon Busch, M., Sánchez-Martínez, J. M., Rodríguez-Martínez, A., Montañez-Enríquez, R., \& Martínez-Jaramillo, S. (2022). Growth at Risk: Methodology and Applications in an Open-Source Platform. \textit{Latin American Journal of Central Banking}, 3, 100068.
    \item Pesaran, M. H. (2021). General Diagnostic Tests for Cross-Sectional Dependence in Panels. \textit{Empirical Economics}, 60(1), 13–50.
    \item Prasad, A., et al. (2019). Growth at Risk: Concept and Application in IMF Country Surveillance. \textit{IMF Working Paper} 2019/036.
    \item Reinhart, C. M., \& Rogoff, K. S. (2009). \textit{This Time Is Different: Eight Centuries of Financial Folly}. Princeton University Press.
    \item Reinhart, C. M., \& Rogoff, K. S. (2011). From Financial Crash to Debt Crisis. \textit{American Economic Review}, 101(5), 1676–1706.
    \item Remolona, E., Scatigna, M., \& Wu, E. (2008). The Dynamic Pricing of Sovereign Risk in Emerging Markets: Fundamentals and Risk Aversion. \textit{Journal of Fixed Income}, 17(4), 57–71.
    \item Segoviano, M., \& Goodhart, C. (2009). Banking Stability Measures. \textit{IMF Working Paper} 09/4.
    \item Stein, J. C. (2012). Monetary Policy as Financial Stability Regulation. \textit{Quarterly Journal of Economics}, 127(1), 57–95.
    \item Uribe, M., \& Yue, V. Z. (2006). Country Spreads and Emerging Countries: Who Drives Whom? \textit{Journal of International Economics}, 69(1), 6–36.
    \item Vasicek, O. (1997). The Loan Loss Distribution. \textit{KMV Corporation}.
    \item Wooldridge, J. M. (2010). \textit{Econometric Analysis of Cross Section and Panel Data} (2nd ed.). MIT Press.
\end{itemize}

\end{document}
